%%%%%%%%%%  scan idx 249  |  folio 236  %%%%%%%%%%
\SGAPrintPage{236}
\sgaStableTarget{sga6.e03.cor.1-12}\textbf{Corollary 1.12.} \uline{In the situation of} (1.11), \uline{if} \(g\) \uline{is pseudo-coherent}, \uline{we have}, \uline{for} \(r\in\mathbb Z\cup\{-\infty\}\),
\[
  D^{-}(f)_{r\text{-}\coh}=D^{-}(h)_{r\text{-}\coh};
\]
\uline{in particular} (\uline{taking} \(X=Y\)), \uline{the notions of} \(r\)-\uline{pseudo-coherence over} \(Y\), ``\uline{in the absolute sense}'' \uline{or rel.} \uline{to} \(g\), \uline{coincide}.

\sgaStableTarget{sga6.e03.cor.1-13}\textbf{Corollary 1.13.} \uline{In the situation of} (1.11), \uline{if} \(f\) \uline{is pseudo-coherent}, \uline{the functor} \(\underline{O}_X\otimes_Y^{\mathrm L}\cdot=Lf^*:D^{-}(Y)\longrightarrow D^{-}(X)\) \uline{induces}, \uline{for} \(r\in\mathbb Z\cup\{-\infty\}\), \uline{a functor}
\[
  Lf^*:D^{-}(g)_{r\text{-}\coh}\longrightarrow D^{-}(h)_{r\text{-}\coh}.
\]

\sgaStableTarget{sga6.e03.cor.1-14}\textbf{Corollary 1.14.} \uline{In the situation of} (1.11), \uline{if} \(f\) \uline{is} \(r\)-\uline{pseudo-coherent} \uline{and} \(g\) \uline{is} \(s\)-\uline{pseudo-coherent} (\(r,s\in\mathbb Z\cup\{-\infty\}\)), \uline{then} \(h\) \uline{is} \(\operatorname{Sup}(r,s)\)-\uline{pseudo-coherent}. \uline{If} \(h\) \uline{is} \(r\)-\uline{pseudo-coherent} \uline{and} \(g\) \uline{is} \((r+1)\)-\uline{pseudo-coherent}, \(f\) \uline{is} \(r\)-\uline{pseudo-coherent}.

\sgaStableTarget{sga6.e03.rem.1-14-1}\textbf{Remark 1.14.1.} In particular, for \(r=s=-1\), we recover (cf. (1.2.1 c))) assertions (ii) and (v) of (EGA IV 1.4.3).

\sgaStableTarget{sga6.e03.sec.2}\section*{2. The finiteness theorem.}

\sgaStableTarget{sga6.e03.conj.2-1}\textbf{Conjecture 2.1.} \uline{Let} \(S\) \uline{be a scheme}, \(f:X\to Y\) \uline{be a} proper \(S\)-\uline{morphism of schemes locally of finite type over} \(S\), \(E\in\Ob D^+(X)\). \uline{Then, if} \(E\) \uline{is pseudo-coherent rel.} \uline{to} \(S\) (1.2), \uline{the same holds for} \(Rf_*(E)\).

\sgaStableTarget{sga6.e03.thm.2-2}\textbf{Theorem 2.2.} \uline{The preceding conjecture is true in the following cases}:
\begin{enumerate}[label=(\roman*)]
\item \(S\) \uline{is locally noetherian};
\item \(f\) \uline{is projective}.
\end{enumerate}

%%%%%%%%%%  scan idx 250  |  folio 237  %%%%%%%%%%
\SGAPrintPage{237}
\uline{Proof in case} (i). Since \(S\) is locally noetherian, the projection of \(X\) (resp. \(Y\)) onto \(S\) is a pseudo-coherent morphism (1.3), and the notions of pseudo-coherence over \(X\) (resp. \(Y\)) and relative to \(S\) coincide (1.12). On the other hand, since \(X\) (resp. \(Y\)) is locally noetherian, hence \(\underline{O}_X\) (resp. \(\underline{O}_Y\)) coherent, an object \(M\) of \(D(X)\) (resp. \(D(Y)\)) is pseudo-coherent if and only if (I 3.5) the \(H^i(M)\) are coherent and \(M\) is locally in \(D^{-}\). The theorem therefore follows from the following one:

\sgaStableTarget{sga6.e03.thm.2-2-1}\textbf{Theorem 2.2.1.} \uline{Let} \(f:X\to Y\) \uline{be a proper morphism of locally noetherian schemes}, \uline{and let} \(E\in\Ob D^+(X)\). \uline{If the} \(H^i(E)\) \uline{are coherent}, \uline{the same holds for the} \(R^if_*(E)\). \uline{If, moreover,} \(E\) \uline{is locally in} \(D^{-}\), \uline{the same holds for} \(Rf_*(E)\).

\uline{Proof}. The first assertion follows from the spectral sequence
\begin{equation}\tag{$*$}
  R^pf_*(H^q(E)) \Longrightarrow R^*f_*(E)
\end{equation}
and the finiteness theorem (EGA III 3.2.1). For the second, we may suppose \(Y\) affine; \(X\) is then quasi-compact and separated. Suppose it covered by \(r\) affine open subsets. We know (EGA III 1.4.12) that, under these conditions, \(R^if_*(M)=0\) for every quasi-coherent \(\underline{O}_X\)-Module \(M\) and every \(i>r\). If \(E\) is locally in \(D^{-}\), \(E\) is in \(D^{-}\), since \(X\) is quasi-compact; say \(H^i(E)=0\) for \(i>N\). From the spectral sequence (*) we then deduce that \(R^if_*(E)=0\) for \(i>N+r\), whence the assertion.

\uline{Proof in case} (ii). The question being local on \(Y\), we may suppose \(S\) and \(Y\) affine. We then have a commutative diagram with a Cartesian square
\[
\begin{tikzcd}[column sep=large,row sep=large]
& & P^r_{Y'} \arrow[dr,"p'"] & \\
& P^r_Y \arrow[ur,hook,"j"] \arrow[d,"p"'] & & Y' \arrow[d,"g'"]\\
X \arrow[ur,hook,"i"] \arrow[r,"f"] & Y \arrow[urr,hook,"j'"] \arrow[rr,"g"] & & S{,}
\end{tikzcd}
\]

%%%%%%%%%%  scan idx 251  |  folio 238  %%%%%%%%%%
\SGAPrintPage{238}
where \(i,j,j'\) are closed immersions, \(g'\) a smooth affine morphism, and \(p,p'\) the projections of the standard projective spaces. Suppose \(E\) pseudo-coherent rel. to \(S\). By (1.1 (iii)), this means that \((ji)_*(E)\) is pseudo-coherent. We must show that \(Rf_*(E)\) is pseudo-coherent rel. to \(S\), i.e. (1.1 (iii)) that \(j'_*Rf_*(E)\) is pseudo-coherent. But we have
\[
  j'_*Rf_*(E)\simeq Rp'_*(ji)_*(E),
\]
so the assertion will follow from the following more precise result:

\sgaStableTarget{sga6.e03.thm.2-2-2}\textbf{Theorem 2.2.2.} \uline{Let} \(Y\) \uline{be an affine scheme}, \(f:X=P^r_Y\longrightarrow Y\) \uline{be the projection of a standard projective space}, \(E\in\Ob D^b(X)_{\mathrm{coh}}\), \(a\in\mathbb Z\). \uline{Then} \(Rf_*(E)\) \uline{is pseudo-coherent}. \uline{Moreover}, \uline{if} \(E\) \uline{is acyclic in degree} \(>a\), \(Rf_*(E)\) \uline{is acyclic in degree} \(>a+r\), \uline{and there exists an integer} \(N\) \uline{such that}, \uline{for every} \(n\geq N\), \(Rf_*(E(n))\) (\uline{where} \(E(n)=E\otimes\underline{O}_X(n)\)) \uline{is acyclic in degree} \(>a\).

\uline{Proof}. Since \(Y\) is affine, \(X\) is quasi-compact and separated, so (II 2.2.9 a)) \(E\) is isomorphic, in \(D(X)\), to a complex \(F\) bounded above in degree such that, for every \(i\), \(F^i\in\Ob \underline{L}(X)\), i.e. (loc. cit.) is a finite direct sum of sheaves of the type \(\underline{O}_X(-n)\) \((n>0)\). Moreover, if \(E\) is acyclic in degree \(>a\), we may choose \(F\) so that \(F^i=0\) for \(i>a\): this follows formally from the preceding assertion, taking account of (I 2.10 a) and II 2.2.8 a)). Replacing \(E\) by \(F\), we may therefore suppose that \(E^i\in\Ob \underline{L}(X)\) for every \(i\), and \(E^i=0\) for \(i>a\). Since \(X\) is covered by \((r+1)\) affine open subsets, we have, by (EGA III 1.4.10) and the hypercohomology spectral sequence, \(R^qf_*(E)=0\) for \(q>a+r\). Let \(t\) be an integer \(\leq a\), and put \(E'=(0\to E^t\to\ldots\to E^a\to0)\). We have an exact sequence
\[
  0\longrightarrow E'\longrightarrow E\longrightarrow E''\longrightarrow0,
\]

%%%%%%%%%%  scan idx 252  |  folio 239  %%%%%%%%%%
\SGAPrintPage{239}
where \(E''\) satisfies the same conditions as \(E\), but with \(a\) replaced by \(t-1\). Applying \(Rf_*\), we obtain a distinguished triangle
\begin{equation}\tag{$*$}
\begin{tikzcd}
& Rf_*(E'') \arrow[dl,hook] & \\
Rf_*(E') \arrow[rr] & & Rf_*(E) \arrow[ul] & .
\end{tikzcd}
\end{equation}
Since \(t\) can be chosen arbitrarily far in the negative direction, it suffices, for the first assertion, to show that \(Rf_*(E)\) is \((t+r)\)-pseudo-coherent. Now \(Rf_*(E'')\) is acyclic in degree \(\geq t+r\), hence \((t+r)\)-pseudo-coherent. Thus, by (I 2.6), the result will follow once we prove that \(Rf_*(E')\) is perfect. For this, we reduce by dévissage (I 4.10) to the case where \(E'\) is concentrated in degree \(0\) and equal to \(\underline{O}_X(-n)\), \(n>0\). Then, by (EGA III 2.1.16), \(R^rf_*(\underline{O}_X(-n))\) is locally free of finite type, and we have
\[
  Rf_*(\underline{O}_X(-n))\simeq R^rf_*(\underline{O}_X(-n))[-r].
\]
Thus \(Rf_*(\underline{O}_X(-n))\) is perfect, and we deduce the first assertion of (2.2.2). We have already seen that \(Rf_*(E)\) is acyclic in degree \(>a+r\). It remains to prove the last assertion. For \(t<-r+a\), \(Rf_*(E'')\) is acyclic in degree \(>a\), so, taking account of the triangle \((*)\), we are reduced to proving the result for \(E'\) in place of \(E\). Since \(E'\) is bounded in degree and its components are objects of \(\underline{L}(X)\), there exists an integer \(N\) such that, for \(n\geq N\), the components of \(E'(n)\) are finite direct sums of sheaves of the type \(\underline{O}_X(p)\), with \(p\geq0\). Then, by (EGA III 2.1.16) and the hypercohomology spectral sequence, \(Rf_*(E'(n))\) is acyclic in degree \(>a\) for every \(n\geq N\), which completes the proof of (2.2.2), and hence of (2.2).

We have a ``classical'' complement to (2.2 (ii)), which we shall give for later reference:

%%%%%%%%%%  scan idx 253  |  folio 240  %%%%%%%%%%
\SGAPrintPage{240}
\sgaStableTarget{sga6.e03.cor.2-3}\textbf{Corollary 2.3} (EGA III 2.2.1). \uline{Let} \(S\) \uline{be a quasi-compact scheme}, \(f:X\to Y\) \uline{be a projective} \(S\)-\uline{morphism of schemes of finite type over} \(S\), \(L=i^*\underline{O}_P(1)\) \uline{be the invertible sheaf defined by an embedding} \(i:X\to P=P(F)\), \uline{where} \(F\) \uline{is} \uline{a quasi-coherent} \(\underline{O}_Y\)-\uline{Module of finite type}. \uline{For} \(M\in\Ob D(X)\) \uline{and} \(n\in\mathbb Z\), \uline{put} \(M(n)=M\otimes^{\mathrm L} L^{\otimes n}\). \uline{Then}, \uline{if} \(E\in\Ob D^+(X)\) \uline{is pseudo-coherent rel.} \uline{to} \(S\) \uline{and} \uline{acyclic in degree} \(>a\) \((a\in\mathbb Z)\), \uline{there exists an integer} \(N\) \uline{such that}, \uline{for every} \(n\geq N\), \(Rf_*E(n)\) \uline{satisfies the same conditions}.

\uline{Proof}. We reduce trivially, as in the proof of (2.2 (ii)), to the situation of (2.2.2): \(Y=S\), \(f:X=P^r_Y\to Y\) the canonical projection, \(L=\underline{O}_X(1)\) the canonical sheaf, \(E\in\Ob D^b(X)\) a pseudo-coherent complex acyclic in degree \(>a\), and the result follows by (2.2.2).

\sgaStableTarget{sga6.e03.cor.2-4}\textbf{Corollary 2.4.} \uline{Let} \(S\) \uline{be a quasi-compact scheme}, \uline{and} \(f:X\to Y\) \uline{be a proper} \(S\)-\uline{morphism of schemes of finite type over} \(S\). \uline{If} \(S\) \uline{is noetherian or} \(f\) \uline{is projective}, \uline{the functor} \(Rf_*:D^+(X)\longrightarrow D^+(Y)\) \uline{induces a functor}
\[
  Rf_*:D^b(X)_{S\text{-}\coh}\longrightarrow D^b(Y)_{S\text{-}\coh}
\]

\uline{Proof}. A consequence of (2.2) and the fact that there is then (EGA III 1.4.12 and IV 1.7.21) an integer \(N\) such that \(R^if_*(E)=0\) for every \(i>N\) and every quasi-coherent \(\underline{O}_X\)-Module \(E\).

\sgaStableTarget{sga6.e03.cor.2-5}\textbf{Corollary 2.5.} \uline{Let} \(f:X\to Y\) \uline{be a proper and pseudo-coherent morphism} (1.2), \uline{and} \(E\in\Ob D^+(X)\). \uline{Suppose} \(Y\) \uline{locally noetherian or} \(f\) \uline{projective}. \uline{Then}, \uline{if} \(E\) \uline{is pseudo-coherent} (\uline{in the absolute sense}), \uline{the same holds for} \(Rf_*(E)\). \uline{If} \(Y\) \uline{is quasi-compact}, \uline{the functor} \(Rf_*\) \uline{induces a functor}
\[
  Rf_*:D^b(X)_{\mathrm{coh}}\longrightarrow D^b(Y)_{\mathrm{coh}}.
\]

